% Add this inside your master.tex preamble file
\usepackage{subfiles}
\usepackage{listings}
\usepackage{booktabs}

\begin{document}

\usepackage{tikz}
\usetikzlibrary{shapes.geometric, arrows.meta, positioning, calc}

% Add this block inside your main text section body
\begin{figure*}[t]
    \centering
    \begin{subfigure}{\linewidth}
    \centering
    \begin{tikzpicture}[
        node distance = 1.2cm and 1.5cm,
        box/.style = {rectangle, draw, fill=blue!5, text width=4.0cm, align=center, minimum height=0.8cm, rounded corners=2pt, font=\small},
        qbox/.style = {rectangle, draw, fill=green!5, text width=2.8cm, align=center, minimum height=1.0cm, font=\small\bfseries},
        wbox/.style = {rectangle, draw, fill=orange!5, text width=3.8cm, align=center, minimum height=0.8cm, rounded corners=2pt, font=\small},
        arrow/.style = {-{Stealth[scale=1.2]}, thick},
        darrow/.style = {-{Stealth[scale=1.2]}, thick, dashed, draw=red!80}
    ]

        % --- BOUNDING BOXES FOR CORES ---
        \draw[dashed, fill=gray!2] (-0.5, 0.5) rectangle (4.5, -4.8);
        \node[font=\footnotesize\bfseries, anchor=north west] at (-0.4, 0.4) {CORE 0 (Deterministic Safety Loop)};

        \draw[dashed, fill=gray!2] (10.5, 0.5) rectangle (15.5, -4.8);
        \node[font=\footnotesize\bfseries, anchor=north east] at (15.4, 0.4) {CORE 1 (Utility \& Dispatch Engine)};

        % --- CORE 0 NODES ---
        \node (c0_start) [box] at (2, -0.2) {Iterate 100\,ms Execution Loop};
        \node (c0_read)  [box, below=of c0_start] {Fetch Raw 16-Bit I2C ADC Metrics};
        \node (c0_math)  [box, below=of c0_read] {Calculate $O_2\%$ and $N_2$ Displacement};
        \node (c0_push)  [box, below=of c0_math] {Push Telemetry Packet to Ring Buffer};

        % --- SHARED QUEUE NODE ---
        \node (queue) [qbox] at (7.5, -3.8) {LOCK-FREE SPSC\\RING BUFFER\\(64 Slots)};

        % --- CORE 1 NODES ---
        \node (c1_pop)   [box, fill=orange!5] at (13, -3.8) {Pop Chronological Telemetry Packet};
        \node (c1_check) [box, fill=orange!5, above=of c1_pop] {Evaluate \texttt{requires\_network\_dispatch}};
        \node (c1_latch) [box, fill=orange!5, above=of c1_check] {Verify Cooling Timer: \\ $(t_{\text{current}} - t_{\text{last}}) \ge T_{\text{cooldown}}$};
        
        % --- BACKGROUND THREAD WORKER ---
        \node (worker)   [wbox, right=of c1_latch, xshift=2.5cm, fill=red!5, draw=red] {Asynchronous Network Worker\\(Detached Outbound Call Task)};

        % --- FLOW LINES CORES ---
        \draw [arrow] (c0_start) -- (c0_read);
        \draw [arrow] (c0_read) -- (c0_math);
        \draw [arrow] (c0_math) -- (c0_push);
        \draw [arrow] (c0_push.west) -- ++(-0.4,0) |- (c0_start.west);

        \draw [arrow] (c1_pop) -- (c1_check);
        \draw [arrow] (c1_check) -- (c1_latch);
        \draw [arrow] (c1_latch.east) -- ++(0.4,0) |- (c1_pop.east);

        % --- INTER-CORE PIPELINE TRACKS ---
        \draw [arrow, draw=blue, line width=1.2pt] (c0_push.east) -- (queue.west);
        \draw [arrow, draw=blue, line width=1.2pt] (queue.east) -- (c1_pop.west);
        
        % --- WORKER SPAWN INTERFACE ---
        \draw [darrow] (c1_latch.east) -- node[above, font=\footnotesize, text=black] {Spawn \& Detach} (worker.west);

    \end{tikzpicture}
    \end{subfigure}
    \caption{System block overview of the asynchronous inter-core communication architecture. The deterministic tracking pipeline on Core 0 transfers metrics safely to the utility processing loops on Core 1 using a lock-free Single-Producer Single-Consumer (SPSC) queue layout, shielding life-safety tasks from background telephony network handshakes.}
    \label{fig:ieee_wide_flow}
\end{figure*}


\section{Nomenclature}
\label{sec:nomenclature}
% Configure a tight, compact list layout to match single-column width constraints
\begin{list}{}{\setlength{\leftmargin}{1.5cm}\setlength{\labelwidth}{1.2cm}\setlength{\itemsep}{2pt}}

    \item[\boldmath$V_{\text{room}}$] 
    Total fixed volume of the enclosed containment structure, in cubic meters ($\text{m}^3$).
    
    \item[\boldmath$O_{2,\text{baseline}}$] 
    Nominal baseline concentration of atmospheric oxygen under normal ambient conditions ($21.0\%$).
    
    \item[\boldmath$O_{2,\text{measured}}$] 
    Real-time oxygen concentration value computed by the sensor array, in percent ($\%$).
    
    \item[\boldmath$\Delta\Phi_{O_2}$] 
    Dimensionless fractional coefficient representing the net volumetric displacement of oxygen.
    
    \item[\boldmath$V_{N_2,\text{released}}$] 
    Calculated absolute volume of target asphyxiant (Nitrogen) gas introduced to the space, in cubic meters ($\text{m}^3$).
    
    \item[\boldmath$t_{\text{current}}$] 
    Current systemic time value derived from the high-resolution hardware clock, in seconds ($\text{s}$).
    
    \item[\boldmath$t_{\text{last\_call}}$] 
    Atomic timestamp indicating the exact moment of the most recent outbound telephony dispatch, in seconds ($\text{s}$).
    
    \item[\boldmath$T_{\text{cooldown}}$] 
    Rigid structural lock out time duration enforced between sequential automated notification events, in seconds ($\text{s}$).
    
    \item[\boldmath$\text{Worker}_{\text{Active}}$] 
    Boolean atomic variable monitoring the operational execution state of the asynchronous network thread thread.

\end{list}


% Pull the patch module into your book architecture cleanly
\subfiles{src/emergency_subfile.tex}

\section{Hardware Interface Loop Metrics}
To guarantee system survival under electrical fault conditions, the physical wiring bridge between the monitoring server and the Advanced Axis AX panel utilizes a supervised three-state loop interface via the 55000-820ADV module. The hardware metrics, calculated line impedances, and expected system state transitions are summarized in Table~\ref{tab:ieee_loop_specs}.

\begin{table}[htbp]
    \centering
    % Force the table font to 8pt footnote size to comply with standard IEEE guidelines
    \footnotesize
    \caption{55000-820ADV Supervised Loop Impedance Metrics}
    \label{tab:ieee_loop_specs}
    % Adjust column internal cell padding to fit exactly into 3.5-inch column constraints
    \setlength{\tabcolsep}{3pt}
    \begin{tabular}{l p{1.6cm} l p{2.4cm}}
        \toprule
        \textbf{Physical} & \textbf{Relay} & \textbf{Target Loop} & \textbf{System Action} \\
        \textbf{State} & \textbf{Position} & \textbf{Impedance} & \textbf{Response} \\
        \midrule
        Standby & Normally Closed & $\approx 6.67\text{ k}\Omega$ & Normal Baseline Log Monitoring \\
        \addlinespace[2pt]
        Emergency & Open Circuit & $\approx 20.00\text{ k}\Omega$ & Trip Relay + Outbound VoIP Call \\
        \addlinespace[2pt]
        Line Shear & Severed Wire & $\infty\ \Omega$ (Open) & Master FACP Yellow Trouble LED \\
        \addlinespace[2pt]
        Line Pinch & Shorted Wire & $\le 2.00\ \Omega$ & Reject Trigger, Flag Error State \\
        \bottomrule
    \end{tabular}
\end{table}

\end{document}
